\section{Situaci\'on del sector paltero en Per\'u y La Libertad}

\subsection{Producci\'on y exportaci\'on peruana}

El sector paltero peruano ha experimentado una expansi\'on sostenida en superficie, producci\'on y exportaciones durante la \'ultima d\'ecada \citep{SierraSelva2020, VitteryFlores2023}. La Tabla~\ref{tab:superficie_rendimiento} muestra el incremento de \'area cosechada, volumen producido y rendimiento promedio entre 2015 y 2019, reflejando tecnificaci\'on y maduraci\'on de huertos. Paralelamente, las exportaciones crecieron en valor y diversificaci\'on de destinos, como indican las Tablas~\ref{tab:destinos_export} y~\ref{tab:export_continente}.

\begin{table}[H]
\centering
\caption{Per\'u: Superficie cosechada, producci\'on y rendimiento de palta (2015--2019)}
\label{tab:superficie_rendimiento}
\small
\begin{tabular}{lrrr}
\toprule
\textbf{A\~no} & \textbf{Superficie (ha)} & \textbf{Producci\'on (miles TM)} & \textbf{Rendimiento (TM/ha)} \\
\midrule
2015 & 28.450 & 445 & 15.6 \\
2016 & 31.200 & 498 & 16.0 \\
2017 & 34.800 & 562 & 16.1 \\
2018 & 38.500 & 628 & 16.3 \\
2019 & 42.100 & 695 & 16.5 \\
\bottomrule
\end{tabular}
\parbox{0.94\textwidth}{\small \textit{Nota.} Elaboraci\'on propia con base en \citep{SierraSelva2020}, Cuadro No.~24.}
\end{table}

\begin{table}[H]
\centering
\caption{Per\'u: Exportaci\'on de palta, principales pa\'ises de destino (2015--2019)}
\label{tab:destinos_export}
\small
\begin{tabular}{lrrrrr}
\toprule
\textbf{Destino} & \textbf{2015} & \textbf{2016} & \textbf{2017} & \textbf{2018} & \textbf{2019} \\
\midrule
Pa\'ises Bajos & 145 & 178 & 210 & 265 & 310 \\
Estados Unidos &  98 & 125 & 155 & 198 & 245 \\
Espa\~na        &  72 &  85 &  98 & 112 & 128 \\
Reino Unido    &  55 &  62 &  70 &  82 &  95 \\
Chile          &  42 &  48 &  52 &  58 &  62 \\
\bottomrule
\end{tabular}
\parbox{0.94\textwidth}{\small \textit{Nota.} Elaboraci\'on propia con base en \citep{SierraSelva2020}, Cuadro No.~28. Valores en millones US\$.}
\end{table}

\begin{table}[H]
\centering
\caption{Per\'u: Exportaciones de palta por continente (2015--2019)}
\label{tab:export_continente}
\small
\begin{tabular}{lrrrrr}
\toprule
\textbf{Continente} & \textbf{2015} & \textbf{2016} & \textbf{2017} & \textbf{2018} & \textbf{2019} \\
\midrule
Europa        & 285 & 340 & 410 & 495 & 580 \\
Am\'erica del Norte & 105 & 132 & 165 & 210 & 258 \\
Am\'erica del Sur  &  48 &  55 &  62 &  70 &  78 \\
Asia          &  85 & 102 & 125 & 148 & 175 \\
\bottomrule
\end{tabular}
\parbox{0.94\textwidth}{\small \textit{Nota.} Elaboraci\'on propia con base en \citep{SierraSelva2020}, Cuadro No.~31. Valores en millones US\$.}
\end{table}


\citet{PromPeru2026} reporta que las exportaciones peruanas de palta fresca superaron los 1.360 millones USD en 2025, con vol\'umenes cercanos a 770 mil toneladas. \citet{MIDAGRI2024} confirma la din\'amica expansiva incluso en periodos de ajuste clim\'atico. La Figura~\ref{fig:evolucion_peru} sintetiza la relaci\'on entre producci\'on y exportaci\'on en el periodo 2015--2019, periodo base para comparaciones estructurales del sector.

\incluirfigura{figuras/figura3\_4.jpg}{0.88\linewidth}{Evoluci\'on de producci\'on y exportaci\'on de paltas, periodo 2015--2019, Per\'u}{fig:evolucion_peru}

La concentraci\'on europea y norteamericana de destinos implica exposici\'on a est\'andares fitosanitarios exigentes, volatilidad log\'istica y competencia con M\'exico y Chile. \citet{ComexPeru2025} enfatiza la necesidad de diversificar mercados sin descuidar la profundizaci\'on en Pa\'ises Bajos y Estados Unidos, principales puertas de entrada actuales.

\subsection{Participaci\'on de La Libertad}

La Libertad concentra el principal n\'ucleo productivo de palta Hass en el Per\'u, con participaciones que oscilan entre 32\% y 36\% del volumen nacional seg\'un campa\~na y fuente estad\'istica \citep{VitteryFlores2023, SierraSelva2020}. La Tabla~\ref{tab:produccion_region} detalla la evoluci\'on de producci\'on por regi\'on, donde La Libertad lidera ampliamente sobre Lambayeque, Piura y otras regiones costeras.

\begin{table}[H]
\centering
\caption{Per\'u: Producci\'on de palta por regi\'on (2015--2019)}
\label{tab:produccion_region}
\small
\begin{tabular}{lrrrrr}
\toprule
\textbf{Regi\'on} & \textbf{2015} & \textbf{2016} & \textbf{2017} & \textbf{2018} & \textbf{2019} \\
\midrule
La Libertad  & 185 & 210 & 245 & 280 & 310 \\
Lambayeque   &  95 & 105 & 118 & 130 & 142 \\
Piura        &  72 &  80 &  88 &  95 & 105 \\
Jun\'in       &  45 &  48 &  52 &  55 &  58 \\
Ancash       &  38 &  42 &  46 &  50 &  54 \\
\bottomrule
\end{tabular}
\parbox{0.94\textwidth}{\small \textit{Nota.} Elaboraci\'on propia con base en \citep{SierraSelva2020}, Cuadro No.~25. Valores en miles de toneladas m\'etricas.}
\end{table}


\begin{table}[H]
\centering
\caption{Per\'u: Valor Bruto de la Producci\'on Agropecuaria por regi\'on (2023--2024)}
\label{tab:vbp_region}
\small
\begin{tabular}{lrr}
\toprule
\textbf{Regi\'on} & \textbf{VBP 2023 (millones S/)} & \textbf{VBP 2024 (millones S/)} \\
\midrule
La Libertad  & 8.245 & 8.620 \\
Lambayeque   & 4.180 & 4.350 \\
Piura        & 3.920 & 4.105 \\
Ica          & 6.540 & 6.890 \\
Lima         & 2.870 & 2.950 \\
\bottomrule
\end{tabular}
\parbox{0.94\textwidth}{\small \textit{Nota.} Elaboraci\'on propia con base en \citep{MIDAGRI2025}, Cuadro C.3.}
\end{table}


La Tabla~\ref{tab:vbp_region} sit\'ua el valor bruto de producci\'on agropecuaria regional, evidenciando el peso econ\'omico de La Libertad dentro del agro nacional \citep{MIDAGRI2025}. Valles como Vir\'u, Moche, Chicama y Chao han impulsado la expansi\'on mediante proyectos de riego, inversi\'on extranjera y experiencia acumulada en agroexportaci\'on de ar\'andano, uva y esp\'arrago.

La ventana comercial de La Libertad se alinea con la costa norte peruana (marzo--agosto), sincronizada con picos de despacho en mayo y junio. Esta concentraci\'on temporal puede generar cuellos de botella en plantas de empaque y presi\'on sobre precios cuando m\'ultiples fundos cosechan simult\'aneamente. La Figura~\ref{fig:destinos_peru} muestra los principales pa\'ises de destino de paltas peruanas en el periodo 2015--2019.

\incluirfigura{figuras/figura3\_5.jpg}{0.88\linewidth}{Principales pa\'ises de destino de paltas peruanas (2015--2019)}{fig:destinos_peru}

\subsection{Problemas de planificaci\'on de campa\~nas}

Los productores de La Libertad enfrentan dificultades para planificar campa\~nas ante incertidumbre de precio, variabilidad de calibres demandados y restricciones log\'isticas \citep{VargasCordova2025, Capunay2025}. La planificaci\'on empresarial de exportadores ---programaci\'on naviera, contratos, estandarizaci\'on de calidad--- no siempre se traduce en informaci\'on oportuna para peque\~nos y medianos productores que venden a intermediarios o participan parcialmente en esquemas integrados.

\citet{VargasCordova2025} identifica en cadenas de valor peruanas asimetr\'ias de informaci\'on y poder de negociaci\'on que perjudican al productor primario. La Figura~\ref{fig:cadena_torobamba} representa eslabones y actores en un valle productivo, ilustrando la complejidad de coordinar decisiones desde la huerta hasta el consumidor final.

\incluirfigura{figuras/figura7\_1.jpg}{0.90\linewidth}{Eslabones y actores de la cadena de valor de palta en el valle Torobamba}{fig:cadena_torobamba}

La ausencia de pron\'osticos accesibles dificulta decidir fechas de cosecha, inversiones en fertilizaci\'on tard\'ia, contrataci\'on de mano de obra y negociaci\'on de precios spot. Cuando la campa\~na se planifica solo con base en el a\~no anterior, se subestima el riesgo de a\~nos at\'ipicos con sobreoferta global o contracci\'on de demanda en Europa.

\subsection{Dependencia de precios hist\'oricos y experiencia emp\'irica}

Muchos productores basan decisiones de siembra y manejo agron\'omico en precios recientes, relatos de vecinos o recomendaciones de intermediarios, m\'etodos que no capturan tendencias, estacionalidad ni choques ex\'ogenos \citep{ReisFilho2022}. \citet{VitteryFlores2023} advierten que mirar exclusivamente el \'ultimo precio FOB puede inducir prociclicidad: ampliar siembras cuando los precios son altos y contraer cuando ya es tarde para evitar p\'erdidas estructurales.

La experiencia emp\'irica aporta conocimiento contextual ---calidad de suelos, microclimas, calendarios locales--- pero no sustituye an\'alisis cuantitativo de series temporales. \citet{Theofilou2025} demuestra en una revisi\'on sistem\'atica que las t\'ecnicas predictivas modernas superan en m\'ultiples estudios a m\'etodos basados en promedios m\'oviles simples o extrapolaci\'on ingenua.

En La Libertad, la brecha digital agrava la dependencia de se\~nales informales: acceso limitado a dashboards de mercado, reportes t\'ecnicos o simuladores de rentabilidad. Cerrar esta brecha requiere traducir modelos estad\'isticos y de aprendizaje autom\'atico en interfaces comprensibles, preferentemente m\'oviles, orientadas a decisiones concretas de siembra y manejo de campa\~na.
